\section{Introduction}

Chronic stress is among the most potent environmental risk factors for neuropsychiatric disease, contributing to the pathogenesis of major depressive disorder, generalized anxiety disorder, and post-traumatic stress disorder. In humans, repeated negative thinking (also termed rumination or worry) represents a cognitive correlate of chronic stress that has been prospectively associated with accelerated cognitive decline and with amyloid and tau accumulation in the brain \cite{marchant2020rnt}. This convergence between a modifiable cognitive process and the biological hallmarks of neurodegeneration raises a causal question: does the persistent reactivation of negatively valenced memory representations itself drive downstream cellular and behavioral pathology, or does it merely co-occur with stress-induced changes driven by other routes? Answering this question requires a model in which a specific, experimentally defined negative memory can be reactivated on demand over long timescales, while cellular and behavioral consequences are measured in a controlled manner.

The ventral hippocampus (vHPC), and specifically the ventral CA1 (vCA1) subregion, is a compelling target for such a model. Anatomical, behavioral, and transcriptomic studies converge on vCA1 as a principal integrator of affective and mnemonic information: the ventral hippocampus is functionally dissociable from its dorsal counterpart, with ventral subfields aligning with emotion and stress circuitry rather than purely spatial or cognitive processing \cite{fanselow2010ventral}. Projections from vCA1 to the basolateral amygdala, nucleus accumbens, and medial prefrontal cortex are organized into largely segregated populations that route distinct behaviorally contingent signals to each target, with anxiety-related firing enriched in prefrontal-projecting neurons \cite{ciocchi2015vca1}. These circuit properties place vCA1 in a privileged position to link a learned aversive experience to enduring shifts in affect, vigilance, and context-dependent behavior.

Advances in activity-dependent tagging now make it possible to identify, label, and manipulate the sparse ensembles that encode specific memories. Optogenetic reactivation of hippocampal dentate gyrus cells active during fear conditioning is sufficient to elicit freezing \cite{liu2012engram}, and the valence of a hippocampal contextual memory engram is malleable, such that pairing reactivation with a stimulus of opposite valence can reverse behavioral responses \cite{redondo2014valence}. Conversely, chronic optogenetic reactivation of dentate cells tagged during a positive experience can rescue stress-induced depressive-like phenotypes \cite{ramirez2015positive}. These findings establish engram ensembles as both necessary and sufficient components of memory expression \cite{josselyn2020engram}, but the complementary question---whether repeated activation of a negative engram over months is sufficient to produce stress-like neural and behavioral pathology---has not been systematically tested.

Any such test must also contend with the well-documented age-dependence of stress vulnerability. Epidemiological data from large international surveys show that childhood adversities account for a substantial fraction of first-onset psychiatric disorders, with effects spanning the entire life-course \cite{kessler2010childhood}. The classical glucocorticoid cascade hypothesis proposes that lifelong exposure to glucocorticoids progressively compromises hippocampal function \cite{sapolsky1986glucocorticoid}, and adult hippocampal neurogenesis---an important substrate for stress resilience---declines sharply with age \cite{klempin2007neurogenesis}. Whether a chronically activated negative engram produces indistinguishable phenotypes in younger and older animals, or whether age modulates the behavioral and glial consequences, has direct implications for how we interpret stress-related disorders across the lifespan.

Here we address these questions using TRAP2 mice \cite{denardo2019temporal}, which provide tamoxifen-gated activity-dependent genetic access to the subset of vCA1 pyramidal neurons active during contextual fear conditioning. We express an excitatory chemogenetic receptor (hM3Dq) in this tagged negative engram and drive it chronically for three months using the high-affinity, metabolically stable DREADD actuator deschloroclozapine (DCZ) \cite{nagai2020dcz,roth2016dreadds} delivered at low dose in drinking water. We perform this manipulation in parallel cohorts that begin at 3 months and 11 months of age, so that the endpoint assays at 6 and 14 months isolate age-dependent differences in the response to an otherwise identical chronic engram-activation protocol. We then ask whether chronic negative engram activation is sufficient to induce: (i) anxiety-like behavior in open field and zero maze; (ii) spatial working memory deficits; (iii) impaired fear extinction and enhanced fear generalization to a novel context \cite{maren2013contextual,tovote2015circuits}; and (iv) cellular signatures of neuroinflammation in the form of microglial and astrocytic remodeling---without overt neuronal loss.

\paragraph{Contributions.} This study makes the following contributions:
\begin{itemize}
  \item We establish a three-month chronic chemogenetic paradigm using TRAP2-tagged vCA1 negative engrams combined with low-dose DCZ-in-water DREADD stimulation, and show that it is well tolerated (no loss of NeuN+ neurons in subiculum, vDG, or vCA1; normal weight gain across both age groups).
  \item We demonstrate that chronic negative engram activation is sufficient to increase anxiety-like behavior in both young (6-month) and aged (14-month) mice, as measured by decreased center time in the open field and decreased open-arm time in the zero maze.
  \item We show that chronic negative engram activation produces age-dependent cognitive effects: spatial working memory, indexed by Y-maze spontaneous alternation, is selectively reduced in 14-month hM3Dq mice, while fear extinction is selectively impaired in 6-month hM3Dq mice, and fear generalization to a novel context is increased at both ages.
  \item We provide evidence that chronic negative engram activation remodels hippocampal glia, with age-dependent microglial (Iba1+) expansion and morphological change in 14-month mice and astrocytic (GFAP+) expansion and branch-length increases at both ages---linking a targeted cognitive manipulation to cellular hallmarks typically associated with chronic stress and early neuroinflammation.
\end{itemize}

\section{Related Work and Background}

\paragraph{Engram tagging and causal manipulation of memory.} The operational idea that a specific memory is supported by a sparse, identifiable neural ensemble has been repeatedly validated using activity-dependent promoters combined with optogenetic or chemogenetic effectors. Optogenetic reactivation of dentate gyrus cells active during fear conditioning is sufficient to evoke freezing in a neutral context \cite{liu2012engram}, and the behavioral valence of an engram can be switched by pairing its reactivation with a stimulus of opposite sign \cite{redondo2014valence}. Reactivation of positive-experience engrams can rescue stress-induced depressive-like phenotypes \cite{ramirez2015positive}, and broader reviews have consolidated these findings into a framework in which engrams are the fundamental units of memory \cite{josselyn2020engram}. Relative to prior engram manipulation studies, which typically involve acute or short-term reactivation over hours to days, our work extends reactivation to three continuous months and focuses specifically on the chronic consequences of reactivating an aversive (fear-conditioned) ensemble in vCA1.

\paragraph{Chemogenetic and TRAP2 tools.} Activity-dependent tagging via targeted recombination in active populations (the TRAP/TRAP2 system) exploits the cFos promoter together with tamoxifen-gated CreER to capture ensembles active during a defined temporal window, and has been used to track and manipulate fear-related cortical ensembles over time \cite{denardo2019temporal}. DREADDs, especially the excitatory hM3Dq receptor, have become standard tools for neuronal manipulation because of their reversibility and cell-type specificity \cite{roth2016dreadds}. However, the canonical DREADD agonist clozapine-N-oxide has pharmacokinetic and off-target liabilities, motivating the development of deschloroclozapine (DCZ), a high-affinity, selective, fast-acting actuator that reliably drives hM3Dq-mediated neuronal activation in rodents and non-human primates at doses of 1--3~$\mu$g~kg$^{-1}$ \cite{nagai2020dcz}. Our low-dose DCZ-in-water protocol builds on this pharmacology but introduces chronic, low-intensity dosing to mimic sustained dysregulation rather than the acute, high-amplitude activation that dominates earlier DREADD studies.

\paragraph{Ventral hippocampus in affect, stress, and aging.} The ventral hippocampus is the hippocampal subfield most closely linked to emotional processing, stress reactivity, and affective disorder \cite{fanselow2010ventral}. Single-cell recordings during anxiety-related behaviors show that subsets of vCA1 projection neurons preferentially route anxiety- or reward-related information to prefrontal cortex, nucleus accumbens, or amygdala \cite{ciocchi2015vca1}. At the systems level, childhood adversity exerts lifelong, cross-diagnostic risk for psychopathology \cite{kessler2010childhood}, consistent with the glucocorticoid cascade model in which repeated stress progressively degrades hippocampal function \cite{sapolsky1986glucocorticoid}, and with a decline in adult hippocampal neurogenesis over the lifespan \cite{klempin2007neurogenesis}. Our design directly addresses the age dimension by running an identical chemogenetic protocol in cohorts that reach 6 versus 14 months of age at endpoint.

\paragraph{Fear circuitry, extinction, and generalization.} Fear memory, extinction, and generalization are supported by a hippocampus-amygdala-medial-prefrontal network in which context integration is especially dependent on the hippocampus \cite{maren2013contextual,tovote2015circuits}. Impairments in extinction and overgeneralization of fear to non-threatening contexts are core phenotypes in anxiety and PTSD models. Our extinction and Context~B (generalization) assays use this framework to read out whether chronic reactivation of a vCA1 negative engram produces deficits that mirror human anxiety-spectrum pathology.
